% Created 2025-01-09 jue 14:07
% Intended LaTeX compiler: pdflatex
\documentclass[presentation]{beamer}
\usepackage[utf8]{inputenc}
\usepackage[T1]{fontenc}
\usepackage{graphicx}
\usepackage{longtable}
\usepackage{wrapfig}
\usepackage{rotating}
\usepackage[normalem]{ulem}
\usepackage{amsmath}
\usepackage{amssymb}
\usepackage{capt-of}
\usepackage{hyperref}
\usepackage{minted}
\usetheme{metropolis}
\usepackage{adjustbox}
\usetheme{default}
\author{Tu Nombre}
\date{\today{}}
\title{Herramientas de Integración y Transferencia de Datos}
\hypersetup{
 pdfauthor={Tu Nombre},
 pdftitle={Herramientas de Integración y Transferencia de Datos},
 pdfkeywords={},
 pdfsubject={},
 pdfcreator={Emacs 28.1 (Org mode 9.6.17)}, 
 pdflang={English}}
\begin{document}

\maketitle

\begin{frame}[label={sec:orgb6b4f1b}]{Introducción}
La transferencia eficiente de datos es crucial en Big Data. Herramientas como:
\begin{itemize}
\item \alert{\alert{Apache Sqoop}}: Transferencia de datos estructurados.
\item \alert{\alert{Apache Flume}}: Ingesta de datos no estructurados.
\end{itemize}

---
\end{frame}

\begin{frame}[label={sec:orge5d0d03}]{Apache Sqoop: Transferencia de Datos Estructurados}
\begin{block}{Características Clave}
\begin{enumerate}
\item \alert{\alert{Interoperabilidad con Bases de Datos Relacionales}}:
\begin{itemize}
\item Compatible con MySQL, PostgreSQL, Oracle, SQL Server, DB2.
\end{itemize}
\item \alert{\alert{Transferencia Bidireccional}}:
\begin{itemize}
\item Importar datos desde bases de datos hacia Hadoop.
\item Exportar datos procesados desde Hadoop hacia las bases de datos.
\end{itemize}
\item \alert{\alert{Optimización Automática}}:
\begin{itemize}
\item Particionamiento de datos para optimizar la transferencia.
\end{itemize}
\item \alert{\alert{Compatibilidad con el Ecosistema Hadoop}}:
\begin{itemize}
\item Integración con HDFS, Hive, y HBase.
\end{itemize}
\end{enumerate}
\end{block}

\begin{block}{Ventajas}
\begin{itemize}
\item Simplifica la integración de datos estructurados.
\item Alta escalabilidad y eficiencia.
\item Compatible con múltiples bases de datos.
\end{itemize}
\end{block}

\begin{block}{Desventajas}
\begin{itemize}
\item No adecuado para flujos en tiempo real.
\item Limitado a datos estructurados.
\end{itemize}
\end{block}

\begin{block}{Casos de Uso}
\begin{itemize}
\item Migración de datos históricos hacia Hadoop.
\item Integración de datos empresariales.
\item Exportación de resultados analíticos.
\end{itemize}
\end{block}

\begin{block}{Ejemplo Práctico}
\begin{itemize}
\item \alert{\alert{Importar datos desde MySQL a HDFS}}:
```bash
sqoop import $\backslash$
  --connect jdbc:mysql://localhost/database\textsubscript{name} $\backslash$
  --username user $\backslash$
  --password password $\backslash$
  --table table\textsubscript{name} $\backslash$
  --target-dir /user/hadoop/target\textsubscript{directory}
```
\item \alert{\alert{Exportar datos desde HDFS a MySQL}}:
```bash
sqoop export $\backslash$
  --connect jdbc:mysql://localhost/database\textsubscript{name} $\backslash$
  --username user $\backslash$
  --password password $\backslash$
  --table table\textsubscript{name} $\backslash$
  --export-dir /user/hadoop/source\textsubscript{directory}
```
\end{itemize}

---
\end{block}
\end{frame}

\begin{frame}[label={sec:org130638b}]{Apache Flume: Transferencia de Datos No Estructurados}
\begin{block}{Características Clave}
\begin{enumerate}
\item \alert{\alert{Ingesta Continua}}:
\begin{itemize}
\item Ideal para logs, eventos, y datos en tiempo real.
\end{itemize}
\item \alert{\alert{Arquitectura Flexible}}:
\begin{itemize}
\item Componentes clave:
\begin{itemize}
\item \alert{\alert{Source}}: Captura datos desde logs o eventos.
\item \alert{\alert{Channel}}: Buffer para garantizar la entrega.
\item \alert{\alert{Sink}}: Envía los datos al sistema de almacenamiento final.
\end{itemize}
\end{itemize}
\item \alert{\alert{Alta Escalabilidad}}:
\begin{itemize}
\item Soporta configuraciones distribuidas.
\end{itemize}
\item \alert{\alert{Tolerancia a Fallos}}:
\begin{itemize}
\item Garantiza la entrega en caso de fallos.
\end{itemize}
\end{enumerate}
\end{block}

\begin{block}{Ventajas}
\begin{itemize}
\item Ideal para flujos en tiempo real.
\item Arquitectura robusta y escalable.
\item Compatible con múltiples fuentes y destinos.
\end{itemize}
\end{block}

\begin{block}{Desventajas}
\begin{itemize}
\item No adecuado para datos estructurados.
\item Transferencia unidireccional.
\end{itemize}
\end{block}

\begin{block}{Casos de Uso}
\begin{itemize}
\item Ingesta de logs hacia HDFS.
\item Flujos de datos IoT.
\item Captura de eventos en tiempo real.
\end{itemize}
\end{block}

\begin{block}{Ejemplo Práctico}
\begin{itemize}
\item \alert{\alert{Configurar un flujo de datos de logs a HDFS}}:
\begin{itemize}
\item \alert{\alert{Archivo de Configuración (flume.conf):}}
```properties
agent.sources = source1
agent.channels = channel1
agent.sinks = sink1

agent.sources.source1.type = exec
agent.sources.source1.command = tail -F /var/log/syslog
agent.sources.source1.channels = channel1

agent.channels.channel1.type = memory
agent.channels.channel1.capacity = 1000
agent.channels.channel1.transactionCapacity = 100

agent.sinks.sink1.type = hdfs
agent.sinks.sink1.hdfs.path = hdfs://localhost:8020/logs/
agent.sinks.sink1.hdfs.fileType = DataStream
```
\item \alert{\alert{Ejecutar el Agente de Flume}}:
```bash
flume-ng agent --conf conf --name agent --conf-file flume.conf
```
\end{itemize}
\end{itemize}

---
\end{block}
\end{frame}

\begin{frame}[label={sec:org04dd658}]{Comparación: Apache Sqoop vs Apache Flume}
\begin{LATEX}\centering
\begin{adjustbox}{width=\linewidth}
\begin{tabular}{|l|l|l|}
\hline
\textbf{Aspecto} & \textbf{Apache Sqoop} & \textbf{Apache Flume} \\
\hline
Tipo de Datos    & Datos estructurados (bases de datos). & Datos no estructurados (logs, eventos). \\
Dirección        & Transferencia bidireccional. & Transferencia unidireccional. \\
Velocidad        & Procesos por lotes (batch). & Flujos en tiempo real (streaming). \\
Escalabilidad    & Alta, optimizada para grandes volúmenes. & Alta, optimizada para flujos continuos. \\
Casos de Uso     & Migración de datos entre bases y Hadoop. & Ingesta de logs, eventos y flujos. \\
\hline
\end{tabular}
\end{adjustbox}
\end{LATEX}

---
\end{frame}

\begin{frame}[label={sec:org7dfc674}]{Conclusión}
\begin{itemize}
\item \alert{\alert{Apache Sqoop}}:
\begin{itemize}
\item Ideal para datos estructurados y transferencias bidireccionales.
\end{itemize}
\item \alert{\alert{Apache Flume}}:
\begin{itemize}
\item Diseñado para flujos en tiempo real y datos no estructurados.
\end{itemize}
\item La elección depende del tipo de datos y la necesidad específica del proyecto.
\end{itemize}
\end{frame}
\end{document}
